\scanpage{58}
% archive foliation: 57
d'un foncteur $E\unsure{\mapsto|\to}\varphi_{T,a}(E)$ des $\pi_1(S,a)$-ensembles finis dans lui-même, foncteur qui est exact à gauche. Je dis qu'un tel foncteur est nécessairement de la forme $\uline{\operatorname{Hom}}_{\pi_1}(A,E)$, où $A$ est un ens, dépendant \deleted{\unsure{et elle}} fonctoriellement du foncteur $\varphi_{T,\unsure{a|c}}$ qui lui \unsure{donne naissance} [produits transformés en sommes]. On pose $A=A(T)$, \deleted{\unsure{le}} et on désigne par $f_*^{\mathrm{ét}}(T)$ la variété \deleted{\unsure{associée}} \inserted{étale} de $S$ associée.
\margin{à détailler, et \unsure{essayer de coiffer par} un th. plus général. Il y a un pt de \unsure{compatibilité} à établir, qui \unsure{résulte} de la formule analogue à $(*)$ ci-dessous, \unsure{si on considère} \uline{\unsure{des ens.}}.}
On a donc un isom de bifoncteurs
\[
(*)\qquad\uline{\operatorname{Hom}}_S(f_*^{\mathrm{ét}}(T),S')=\deleted{\unsure{f_*}}\uline{\operatorname{Hom}}_S(T,S')
\]
[sur $\operatorname{Rev\ ét}(T)\times\operatorname{Rev\ ét}(S)$]

\deleted{Il reste à montrer que tt $T$, le foncteur en $S'$ \unsure{du|des} deuxième \unsure{mbs|mb} est représentable.}
Prenant \deleted{\unsure{les sections}, $S'=T$, et} \deleted{\unsure{les deux foncteurs}} les sections des deux côtés, on trouve
\[
\operatorname{Hom}_S(\unsure{f_\alpha^{\mathrm{ét}}|f_*^{\mathrm{ét}}}(T),S')\simeq\operatorname{Hom}_S(T,S')
\]
d'où, en faisant $S'=\unsure{f_\alpha^{\mathrm{ét}}|f_*^{\mathrm{ét}}}(T)$, \unsure{on tire|on trouve}
